\section{A review of a classical PCP theorem}\label{sec:classical-pcp}

We begin \cref{part:neexp} by reviewing the following classical PCP theorem:
\begin{equation}
\succinct \in \mathrm{PCP}[n, \poly(n)].
\end{equation}
This implies, by standard reduction, that $\succinct \in \mathrm{MIP}$,
which is the main result of~\cite{BFL91}.
Reviewing this serves two purposes:
(i) our $\mip^*$ protocols are inspired by this PCP construction,
and (ii) their correctness is actually shown by reduction to the correctness of this PCP construction (\Cref{lem:polynomial-means-good} below).
This section closely follows the excellent course notes of Harsha~\cite{Har10}.

\subsection{The instance}

The input to the verifier is an instance of the $\succinct$ problem,
i.e.\ a circuit~$\calC$  of size~$s$ with $3n+3$ inputs.
We apply the Tseitin transformation to it to produce a formula~$\calF$ with $n' = 3n+3+s$ inputs.
It encodes the $\sat$ formula $\psi := \psi_{\calF}$ on~$N = 2^n$ variables
which contains $(x_i^{b_1} \lor x_j^{b_2} \lor x_k^{b_3})$ 
as a clause
if and only if~$\calF(i, j, k, b_1, b_2, b_3, w) = 1$ for some $w \in \{0, 1\}^s$.

\subsection{Encoding assignments}

Writing $\calS = \{0, 1\}^n$, an assignment to the variables of~$\psi$ is a string $a \in \{0,1\}^{\calS}$, or equivalently a string in $\{0, 1\}^N$.
The first step of the PCP theorem is to take the low-degree encoding of~$a$.
We begin by choosing parameters.

\begin{definition}\label{def:admissible}
Recall that
$N = 2^n$,
$h= 2^{t_1}$,
$q = 2^{t_2}$,
and $m$ are admissible parameters if $t_1 \leq t_2$ and $h^m \geq N$.
We call them \emph{exactly} admissible if the stronger condition $h^m = N = 2^n$ holds.
\end{definition}

We select $n$, $h = 2^{t_1}$, $q = 2^{t_2}$, and $m$ to satisfy our ``rule of thumb" parameter settings (\Cref{eq:parameters}):
\begin{equation*}
	h = \Theta(n), \quad m = \Theta\left(\frac{n}{\log(n)}\right), \quad q = \Theta((n')^{10}).
\end{equation*}
Note that $q$ depends on~$n'$ rather than just~$n$.

It remains to choose~$H$ and~$\pi$.
Our construction requires that the permutation~$\pi$ be efficiently computable,
and so we pick these according to the canonical low-degree encoding (\Cref{def:canonical-low-degree}).
This entails setting $H= H_{t_1, t_2}$.
As for $\pi$, we modify the construction slightly.
This is because the canonical low-degree encoding is designed for strings whose coordinates are indexed by an integer $i \in[n]$,
which must first be converted to binary when computing~$\pi$.
However, the coordinates of our strings $a \in \{0, 1\}^{\calS}$ are indexed by elements of $\calS = \{0, 1\}^n$,
which are already written in binary, allowing us to skip the conversion.
Hence, within this section, we define $\pi:=\pi_{n, t_1, t_2} :\calS \rightarrow H^m$ by setting
\begin{equation*}
\pi(b_1, \ldots, b_n) := \sigma_{n, t_1, t_2}(b_1, \ldots, b_n)= (\sigma(b_1, \ldots, b_{t_1}), \sigma(b_{t_1+1}, \ldots, b_{2t}), \ldots, \sigma(b_{n-t_1+1}, \ldots, b_n)),
\end{equation*}
where $\sigma:=\sigma_{t_1, t_2}$.
Given these parameters,
an assignment~$a$ is encoded as a degree-$O(mh)$ polynomial $g_a:\F_q^m \rightarrow \F_q$.

The crucial property of~$\pi$ that we will need later is that it has an efficiently-computable, low-degree inverse.
We will show this here.
To do so, we begin by recalling the notation $\indicator{H}{x}{y}$ for the indicator function of~$x \in H$ over~$H$:
\begin{equation*}
\indicator{H}{x}{y} = \frac{\prod_{b \neq x}(y - b)}{\prod_{b \neq x}(x-b)},
\end{equation*}
where~$b$ ranges over~$H$.
This is a degree-$h$ polynomial which can be computed in time $\poly(h, q)$.
\begin{definition}
Let $N= 2^n$, $h = 2^{t_1}$, $q= 2^{t_2}$, and $m$ be exactly admissible parameters.
Set $H = H_{t_1, t_2}$, $\sigma = \sigma_{t_1, t_2}$, and $\pi = \pi_{n, t_1, t_2}$.
Consider the function $\mu:= \mu_{t_1, t_2}:\F_q \rightarrow \F_q^{t_1}$ whose $i$-th component is defined as
\begin{equation*}
\mu_i(y) = \sum_{x \in H : \tr[e_i \cdot x] = 1} \indicator{H}{x}{y}.
\end{equation*}
Let $y = b_1 \cdot e_1 + \cdots + b_{t_1} \cdot e_{t_1}$ be an element of $H$.
Then $\mu_i(y) = b_i$, and so $\mu(y) = (b_1, \ldots, b_{t_1})$.
This means that $\mu(\sigma(b_1, \ldots, b_{t_1})) = (b_1, \ldots, b_{t_1})$.
As a result, if we define the function $\nu := \nu_{n, t_1, t_2} : \F_q^m \rightarrow \calS_n$ to be
\begin{equation*}
\nu(x_1, \ldots, x_m) := (\mu(x_1), \ldots, \mu(x_m))
\end{equation*}
then $\nu(\pi(x)) = x$ for any $x \in \calS_n$.
Each component of $\nu$ is the sum of $\frac{h}{2}$ indicator functions,
and is therefore degree-$h$ and computable in time $\poly(h,q)$.
As a result, $\nu$ is computable in time $\poly(n,h,q)$.
\end{definition}


\subsection{Encoding the formula}

Our next step is to produce a similar ``low-degree encoding" for the formula~$\psi$.
This will be a function $g_\psi: \F_q^{m'} \rightarrow \F_q$, for $m' = 3m + 3+s$, with the property that for all $i, j, k \in \calS$, $b_1, b_2, b_3 \in \{0, 1\}$, and $w \in \{0, 1\}^s$,
\begin{equation*}
g_\psi(\pi(i), \pi(j), \pi(k), b_1, b_2, b_3, w) = \calF(i, j, k, b_1, b_2, b_3, w).
\end{equation*}
This can be accomplished by setting $\calS' := \{0, 1\}^{n'}$, viewing $\calF$ as computing a string $a_{\psi} \in \{0, 1\}^{\calS'}$, and setting $g_{\psi}$ to be its low-degree encoding.
However, the verifier in our protocol will be required to evaluate $g_\psi$ on a particular input,
and this seems challenging given that this $g_{\psi}$ is computed by interpolating over an exponential number of points.
Instead, we will produce a $g_\psi$ which we can efficiently evaluate at any point using the fact that we have a succinct formula~$\calF$ representing~$\psi$.

To begin, we convert~$\calF$ into an algebraic formula which operates over $\F_q$-valued inputs using arithmetization (cf.\ \Cref{def:arithmetization}).
Set $\calF_{\mathrm{arith}} := \arith{q}{\calF}$.
This is a  function $\calF_{\mathrm{arith}}:\F_q^{n'}\rightarrow \F_q$
with the property that for any $x \in \{0, 1\}^{n'}$, $\calF_{\mathrm{arith}}(x) = \calF(x)$.
Furthermore, $\calF_{\mathrm{arith}}$ has degree~$O(n')$ and is computable in time $\poly(n', q)$.
We can now define the function $g_\psi$ as follows.
\begin{definition}\label{def:encoded-function}
Let $N=2^n$, $h = 2^{t_1}$, $q= 2^{t_2}$, and $m$ be exactly admissible parameters.
Set $\nu:=\nu_{n, t_1, t_2}$.
Let $\calC$ be a $\succinct$ instance whose Tseitin transformation~$\calF$
has $n' = 3n + 3 + s$ inputs and encodes the formula $\psi := \psi_{\calF}$,
and let $\calF_{\mathrm{arith}} = \arith{q}{\calF}$.
Write $m' = 3m + 3 + s$.
Then we define $g_\psi := g_{\psi, n, t_1, t_2} : \F_q^{m'} \rightarrow \F_q$ to be the function
\begin{equation*}
g_\psi(x_1, x_2, x_3, b_1, b_2, b_3, w) := \calF_{\mathrm{arith}}(\nu(x_1), \nu(x_2), \nu(x_3), b_1, b_2, b_3, w).
\end{equation*}
This is degree $h \cdot O(n')$ and can be computed in time $\poly(n',h,q)$.
\end{definition}
For shorthand, we will often write inputs to $g_\psi$ as tuples $(x, b, w) \in \F_q^{3m+3+s}$, where $x = (x_1, x_2, x_3)$ contains three strings in $\F_q^m$
and $b = (b_1, b_2, b_3)$ contains three numbers in $\F_q$.

\subsection{Zero on subcube}

Given a function $g:\F_q^m \rightarrow \F_q$,
we would like to check that it is the low-degree encoding of an assignment which satisfies~$\psi$.
To do so, we define the following function.
\begin{definition}\label{def:formula-function}
Let $N=2^n$, $h = 2^{t_1}$, $q= 2^{t_2}$, and $m$ be exactly admissible parameters.
Let $\calC$ be a $\succinct$ instance whose Tseitin transformation~$\calF$
has $n' = 3n + 3 + s$ inputs and encodes the formula $\psi := \psi_{\calF}$,
and let $g_\psi := g_{\psi, n, t_1, t_2}$.
Set $m' = 3m + 3 + s$.
Then given a function $g:\F_q^m \rightarrow \F_q$, 
we define $\mathrm{sat}_{\psi,g} := \mathrm{sat}_{\psi,g, n, t_1, t_2} : \F_q^{m'} \rightarrow \F_q$ to be the function
\begin{equation*}
\mathrm{sat}_{\psi,g}(x, b,w) := g_{\psi}(x, b, w) \cdot (g(x_1)-b_1) (g(x_2)-b_2) (g(x_3)-b_3).
\end{equation*}
\end{definition}
The crucial property we would like to check is that $\mathrm{sat}_{\psi, g}$ is \emph{zero on the subcube $H_{\mathrm{zero}} := H^{3m} \otimes \{0, 1\}^{3+s}$}.
\begin{proposition}\label{prop:zero-on-subcube}
The function $\mathrm{sat}_{\psi,g}$ is zero on the subcube~$H_{\mathrm{zero}}$ for some $g:\F_q^m\rightarrow \F_q$ if and only if~$\psi$ is satisfiable.
If it \emph{is} satisfiable, $g$ may be taken to be degree-$O(mh)$, in which case $\mathrm{sat}_{\psi,g}$ is  degree-$O(mh + h n')$.
\end{proposition}
Note what \Cref{prop:zero-on-subcube} does \emph{not} say.
It does not say that if $\mathrm{sat}_{\psi,g}$ is zero on the subcube, then~$g$ is the low degree encoding of a satisfying assignment of~$\psi$.
It does not even say that~$g$ must be low-degree.
(Indeed, $g$ might have high degree, as $\mathrm{sat}_{\psi,g}$ only checks~$g$ on those numbers in the range of~$\pi$.)
What it \emph{does} say is that \emph{if} $\psi$ is satisfiable, \emph{then} there exists such a $g$ which is low-degree: the low-degree encoding of a satisfying assignment.
Our strategy, then, will be to verify that that $g$ is low-degree and then use this fact to verify that $\mathrm{sat}_{\psi,g}$ is zero on the subcube.
(We can then ``forget" that~$g$ is low-degree, as it is no longer required for the analysis.)

To verify this that $\mathrm{sat}_{\psi,g}$ is zero on~$H_{\mathrm{zero}}$, we would like it to be encoded so that this is self-evidently true.
This entails expanding $\mathrm{sat}_{\psi,g}$ in a ``basis" of simple polynomials which are zero on the subcube.
To begin, given a subset $S \subseteq \F_q$, define
\begin{equation*}
\mathrm{zero}_S(x) := \prod_{b \in S}(x-b).
\end{equation*}
The following proposition shows how to expand into this ``zero" basis.
\begin{proposition}\label{prop:coefficient-polys}
Let $f:\F_q^{n} \rightarrow \F_q$ be a degree-$d$ polynomial which is zero on the subcube $H = H_1 \otimes \cdots \otimes H_n$. Then there exist degree-$(d-h)$ ``coefficient polynomials" $c_1, \ldots, c_n$ such that
\begin{equation*}
f(x) = \mathrm{zero}_{H,c}(x):= \sum_{i=1}^{n} \mathrm{zero}_H(x_i) \cdot c_i(x).
\end{equation*}
\end{proposition}

For simplicity, we will write $\mathrm{zero}_{H,c}$ instead of $\mathrm{zero}_{H_{\mathrm{zero}},c}$.
We would like our proof to consist of the function~$g$ and the coefficient polynomials $c_1,\ldots, c_{m'}$ so that we may check the equality
$\mathrm{sat}_{\psi,g} \equiv \mathrm{zero}_{H,c}$.
The following lemma shows so long as these functions are low-degree,
we can verify that they are equal, and therefore show that $\psi$ is satisfiable.

\begin{lemma}\label{lem:the-grand-finale}
Let $N=2^n$, $h = 2^{t_1}$, $q= 2^{t_2}$, and $m$ be exactly admissible parameters.
Let $\calC$ be a $\succinct$ instance whose Tseitin transformation~$\calF$
has $n' = 3n + 3 + s$ inputs and encodes the formula $\psi := \psi_{\calF}$,.
Set $m' = 3m + 3 + s$.
Let $g:\F_q^m \rightarrow \F_q$, and set $\mathrm{sat}_{\psi,g} := \mathrm{sat}_{\psi,g, n, t_1, t_2}$.
Let $c_1, \ldots, c_{m'}: \F_q^{m'} \rightarrow \F_q$, set $H_{\mathrm{zero}} = H^{3m} \otimes \{0, 1\}^{3 + s}$, and write $\mathrm{zero}_{H_c} := \mathrm{zero}_{H_{\mathrm{zero}},  c}$.
Suppose that $g$ is degree-$d_1$, and suppose that $c_1, \ldots, c_{m'}$ are degree-$d_2$.
Suppose
\begin{equation*}
\Pr_{\bx \sim \F_q^{m'}}[\mathrm{sat}_{\psi,g}(\bx) =  \mathrm{zero}_{H,c}(\bx)] > \frac{\max\{O(h n') + 3 d_1, h + d_2\}}{q}.
\end{equation*}
Then $\psi$ is satisfiable.
\end{lemma}
\begin{proof}
By \Cref{def:encoded-function},
$\mathrm{sat}_{\psi, g}$ has degree $h \cdot O(n') + 3 d_1$.
In addition, $\mathrm{zero}_{H,c}$ has degree $h + d_2$.
Define $f = \mathrm{sat}_{\psi,g} -  \mathrm{zero}_{H,c}$.
Then~$f$ has degree $\max\{O(h n') + 3 d_1, h + d_2\}$.
By assumption, $f(\bx) = 0$ with probability larger than $\deg(f)/q$ over a uniformly random $\bx \sim \F_q^{m'}$.
By Schwartz-Zippel (\Cref{lem:schwartz-zippel}), this means that $f \equiv 0$, which implies that $\mathrm{sat}_{\psi, g} \equiv \mathrm{zero}_{H,c}$.
But $\mathrm{zero}_{H,c}$ is self-evidently zero on the subcube $H_{\mathrm{zero}}$,
meaning that $\mathrm{sat}_{\psi,g}$ is as well.
By \Cref{prop:zero-on-subcube}, $\psi$ is satisfiable.
\end{proof}

Ensuring that $\mathrm{sat}_{\psi,g}$ is low-degree requires ensuring that~$g$ is low-degree, and this can be accomplished with the low-degree test.
Arguing similarly for $\mathrm{zero}_{H,c}$ requires ensuring that each~$c_i$ is low-degree, and this can be done with the simultaneous plane-versus-point low-degree test (\Cref{thm:simultaneous-raz-safra}).

\subsection{The PCP}\label{sec:the-pcp}

We can now state the contents of our probabilistically checkable proof for the satisfiability of~$\psi$.
It consists of the following four tables.
\begin{enumerate}
\itemsep -.5pt
\item A claimed low-degree polynomial $g:\F_q^m \rightarrow \F_q$.
\item A set of claimed low-degree polynomials $c_1, \ldots, c_{m'} : \F_q^{m'} \rightarrow \F_q$.
\item A ``planes table", containing for each plane~$s$ in $\F_q^m$ a degree-$d$ bivariate polynomial.
\item Another planes table, containing for each plane~$s$ in $\F_q^{m'}$ an $m'$-tuple of degree-$d$ bivariate polynomials.
\end{enumerate}
The verifier works as follows:
first, it performs the low-degree test between $g$ and its planes table.
Second, it performs the simultaneous low-degree test between the $c_i$'s and their plane table.
Both of these use the degree parameter $d = \Theta((n')^2)$,
which is chosen to upper-bound both $\Theta(mh)$ and $\Theta(mh + hn')$.
Finally, it picks a uniformly random $(\bx, \bb, \bw) \in \F_q^{m'}$
and checks the equality $\mathrm{sat}_{\psi,g}(\bx, \bb, \bw) = \mathrm{zero}_{H,c}(\bx, \bb, \bw)$.
It accepts if all the tests accept individually.

When $\psi$ is satisfiable, there is always a proof that makes the verifier accept with probability~$1$.
This entails setting $g$ to be the low-degree encoding of a satisfying assignment,
and setting $c_1, \ldots, c_{m'}$ to be the coefficient polynomials of $\mathrm{sat}_{\psi,g}$.
The following proposition shows that when~$\psi$ is not satisfiable, the verifier always rejects with probability at least~$\tfrac{1}{10}$.
\begin{proposition}
If the verifier accepts with probability at least $9/10$, then $\psi$ is satisfiable.
\end{proposition}
\begin{proof}
If the verifier accepts with probability at least $9/10$,
then each individual test accepts with probability at least $9/10$.
Applying \Cref{thm:raz-safra,thm:simultaneous-raz-safra},
we get degree-$d$ functions $\overline{g}:\F_q^n \rightarrow \F_q$ and $\overline{c}_1, \ldots, \overline{c}_{m'}:\F_q^{m'} \rightarrow \F_q$
such that
\begin{equation*}
\mathrm{dist}(g, \overline{g}) \leq \tfrac{2}{10}, \quad
\mathrm{dist}(c, \overline{c}) \leq \tfrac{2}{10}, \quad
\mathrm{dist}(\mathrm{sat}_{\psi,g}, \mathrm{zero}_{H,c}) \leq \tfrac{1}{10}.
\end{equation*}
(Here, we are assuming that $q$ is a sufficiently large function of~$m$ and~$h$.)
By the union bound, $\mathrm{dist}(\mathrm{sat}_{\psi,g}, \mathrm{sat}_{\psi,\overline{g}}) \leq 3\mathrm{dist}(g, \overline{g})$.
As a result, by the triangle inequality
\begin{align*}
\mathrm{dist}(\mathrm{sat}_{\psi,\overline{g}}, \mathrm{zero}_{H,\overline{c}})
&\leq \mathrm{dist}(\mathrm{sat}_{\psi,\overline{g}}, \mathrm{sat}_{\psi, g}) +
	\mathrm{dist}(\mathrm{sat}_{\psi, g}, \mathrm{zero}_{H,c})+
	\mathrm{dist}(\mathrm{zero}_{H,c} + \mathrm{zero}_{H,\overline{c}})\\
&\leq 3 \mathrm{dist}(g, \overline{g})+
	\mathrm{dist}(\mathrm{sat}_{\psi, g}, \mathrm{zero}_{H,c})+
	\mathrm{dist}(c, \overline{c}) \leq 3 \cdot \tfrac{2}{10} + \tfrac{2}{10} + \tfrac{1}{10} = \tfrac{9}{10}.
\end{align*}
By \Cref{lem:the-grand-finale}, $\psi$ is therefore satisfiable.
\end{proof}

\paragraph{Time and communication complexity.}
\begin{itemize}
\itemsep -.5pt
\item[$\circ$] \textbf{Question length:} The verifier performs two low-degree tests and draws a random point in $\F_q^{m'}$.
			These are of size $\Theta(m \log(q))$, $\Theta(m'\log(q))$, and $\Theta(m'\log(q))$, respectively,
			all of which are $O(n')$ bits.
\item[$\circ$] \textbf{Answer length:} The verifier performs one normal low-degree test, and then a second low-degree test with answer complexity $m'$ times the normal answer complexity. These are of total length $(m' + 1) \cdot d^2 \log(q) = O((n')^9)$. Finally, in the last test, it queries each of $g$ and $c_1, \ldots, c_{m'}$ for a point in $\F_q$, a total communication cost of $(m'+1) \log(q) = O(n')$. In total, the answer length is $\poly(n')$.
\item[$\circ$] \textbf{Runtime:} The verifier runs in time $\poly(n')$. This includes computing $\mathrm{sat}_{\psi,g}(\bx, \bb, \bw)$,
						which requires computing $g_{\psi}(\bx, \bb, \bw)$, taking time $\poly(n', h, q) = \poly(n')$.
\end{itemize}

%%% Local Variables:
%%% mode: latex
%%% TeX-master: "main"
%%% End:
%%%---------- close: little-nexp-protocol

%%%%%%%%%%% jump to big-low-degree-prelims
%%%---------- open: big-low-degree-prelims
%!TEX root = main.tex

